\documentclass[11pt,letterpaper]{article}
\usepackage{style}
\hypersetup{pdftitle={Radin Shayanfar - CV}}

\begin{document}
\datefooter

\header{LLM Coding Agents \textbullet\ Agent Evaluation \& Benchmarking}

\section*{Summary}
AI/ML engineer and researcher at Huawei Canada working on \textbf{agentic AI}, specifically \textbf{LLM coding agents} and their evaluation. I build, audit, and improve the \textbf{benchmarks} and \textbf{harnesses}, and work to make agents' output \textbf{trustworthy}, with papers accepted at \textbf{NeurIPS} 2026. During my M.A.Sc.\ I worked on controlling LLM behaviour, compiling dialogue schemas into \textbf{interpretable guardrails} (CoDial, \textbf{ACL} 2026), making models safer via \textbf{multi-agent debate} (RedDebate, \textbf{ICML} 2026), and applying \textbf{RAG} to misinformation detection (MisLC, \textbf{EMNLP} 2024). I contribute to NVIDIA NeMo Guardrails and the Harbor framework.

\section*{Experience}

\entry{Associate LLM Engineer}{\href{https://www.huawei.com/}{Huawei Canada}}{Kingston, ON, Canada}{Sep.\ 2025 -- Present}
\begin{itemize}
  \item Lifted gold-patch resolution on FeatBench from 37.31\%\,{\small$\pm$\,1.88} to \textbf{99.81\%}\,{\small$\pm$\,0.10} by diagnosing harness, container, and dependency-pinning defects in the original benchmark, and released the fixes as FeatBench-Verified v1.3.
  \item Raised an in-house model's resolve rate on SWE-Bench Verified by \textbf{15+ points} by auditing agent trajectories and modifying the OpenHands scaffold to mitigate the recurring failure modes they exposed, including a hardened string-replacement tool worth $\sim$10 points of that gain, plus evaluation-infrastructure defects.
  \item Designed and built a \textbf{375-task} multi-language agentic coding benchmark in the SWE-Bench family, spanning bug fixes, feature implementation, documentation, refactoring, and testing. Problem statements are drawn from real developer--agent interactions rather than GitHub issue text, closing the gap between how benchmarks phrase tasks and how developers actually ask for them (paper in preparation).
  \item Shipped each task at three levels of problem detail, varying how much the developer reveals to the agent, so the benchmark measures sensitivity to problem specification rather than a single fixed phrasing.
  \item Profiled what five agentic SE benchmarks demand of an agent by running six coding agents across two model families at three scales each; introduced the Spread--Novelty--Centrality profile, showed that category labels such as ``bug fix'' are unreliable proxies for the engineering work a task requires, and surfaced behavioural patterns specific to each family and scale (first-author, preprint).
  \item Isolated recurring model-side failure modes in LiveCodeBench (LCB) agent trajectories and reported them to the SFT data team to inform the next training iteration; separately surfaced harness and evaluation defects that were masking true capability signal.
\end{itemize}

\entry{Big Data Engineer}{\href{https://bale.ai/}{Bale Messenger}}{Tehran, Iran}{Mar.\ 2023 -- Aug.\ 2023}
\begin{itemize}
  \item Developed and maintained production big data pipelines for continuous ingestion and processing on the Bale financial-messaging platform, serving $\sim$20M monthly active users.
  \item Implemented scalable solutions with Apache Spark, Hive, and Kafka, improving the efficiency and scalability of the data warehouse and the real-time accessibility of its data.
\end{itemize}

\entry{Machine Learning Engineer \& Backend Developer}{Sharif Data Analytics Lab \& National Elites Foundation of Iran}{Tehran, Iran}{Dec.\ 2021 -- Sep.\ 2022}
\begin{itemize}
  \item Trained a speaker-verification system (part of an Intelligent Voice Commands Recognition product) leveraging the state-of-the-art \href{https://arxiv.org/abs/2104.01466}{ECAPA-TDNN} model, achieving over 97\% accuracy on a custom 8,400-pair dataset.
  \item Designed and developed a web API using a microservices architecture with Flask and Laravel, streamlining voice-command processing and authentication.
\end{itemize}

\entry{Scientific Team Member}{\href{https://samcode.allamehelli3.ir/staff.html}{SamCode 5 Programming Contest}}{Tehran, Iran}{Sep.\ 2020 -- Feb.\ 2021}
\begin{itemize}
  \item Designed algorithmic problems and comprehensive test cases for a programming contest targeting high-achieving junior-high students.
\end{itemize}

\entry{Freelance Developer}{\href{https://ponisha.ir/profile/radinshayanfar}{Ponisha}}{Tehran, Iran}{2020}
\begin{itemize}
  \item Built Telegram-based applications in Python using Telegram's TDLib, delivering tailored solutions for a range of clients.
\end{itemize}

\subsection*{Teaching Experience}
\begin{itemize}
  \item \textbf{Queen's University} --- Teaching Assistant: Machine Learning \& Deep Learning (Winter 2025, Prof.\ X.\ Zhu); Fundamentals of Information Structure (Fall 2024, Prof.\ T.\ Dean); Object-Oriented Programming (Winter 2024, Prof.\ M.\ Greenspan).
  \item \textbf{Amirkabir University of Technology} --- Teaching Assistant: Multicore Programming (Spring 2022 \& 2023, Prof.\ M.\ Momtazpour); Computational Intelligence (Spring 2022, Prof.\ M.\ Ebadzadeh); Signals and Systems (\textbf{Head TA}, Fall 2021, Prof.\ M.\ Rasti; Spring 2021, Dr.\ A.\ Aghaeeyan); Data Structures and Algorithms (Fall 2021, Prof.\ E.\ Nazerfard); Linear Algebra (Spring 2021, Prof.\ M.\ H.\ Chehreghani).
\end{itemize}

\section*{Education}

\entry{Master of Applied Science, Electrical and Computer Engineering}{Queen's University, Kingston, ON, Canada}{GPA: A+}{Sep.\ 2023 -- Sep.\ 2025}
\begin{itemize}
  \item Thesis: \textit{CoDial: Interpretable Task-Oriented Dialogue Systems Through Dialogue Flow Alignment} (Prof.\ X.\ Zhu). Published at ACL 2026.
  \item Research areas: NLP, dialogue systems, NLP in the legal domain, machine learning.
\end{itemize}

\entry{Bachelor of Science, Computer Engineering}{Amirkabir University of Technology, Tehran, Iran}{GPA: 19.45 / 20}{Sep.\ 2018 -- Apr.\ 2023}
\begin{itemize}
  \item Ranked \textbf{1st} among 156 students.
  \item B.Sc.\ Project: \textit{Generating Handwritten Persian Characters with Generative Models} (Prof.\ A.\ Nickabadi).
  \item Selected coursework: Computational Intelligence (20/20), Artificial Intelligence (20/20), Data Structures and Algorithms (20/20), Information Retrieval (20/20), Programming Languages (20/20), Signals and Systems (20/20), Advanced Programming (20/20), Web Programming (20/20), Multicore Programming (19.8/20), Multimedia Systems (19.5/20), Data Mining (19.5/20), Operating Systems (19.4/20).
\end{itemize}

\entry{High School Diploma, Mathematics and Physics}{Allame Helli 3 High School (NODET), Tehran, Iran}{Average: 19.16 / 20}{2014 -- 2018}

\section*{Publications \& Writing}

\pubentry[Introduces the Spread--Novelty--Centrality profile to characterize what repository-level coding tasks demand, applied to five benchmarks and six agents.]
  {\href{https://arxiv.org/abs/2609.01271}{What Does an Agentic Software Engineering Benchmark Measure? Profiling Task Demands and Agent Behaviour Beyond What Category Labels Reveal}}
  {\textbf{Radin Shayanfar}, Keheliya Gallaba, Ahmed E.\ Hassan.}
  {Preprint}

\pubentry
  {\href{https://arxiv.org/abs/2609.04298}{Harbor Adapters and Harbor-Index: Infrastructure and a Curated Meta-Dataset for Large-Scale Agentic Evaluation}}
  {Lin Shi, Haowei Lin, \ldots, \textbf{Radin Shayanfar}, \ldots, Ludwig Schmidt, Alex G.\ Shaw.}
  {NeurIPS 2026\\(E\&D Track)}

\pubentry[Controls LLM behaviour from a dialogue flow the developer writes, turning it into guardrail code (NVIDIA Colang) so the model's decisions stay traceable and easy to correct. SOTA on standard TOD benchmarks.]
  {\href{https://aclanthology.org/2026.acl-long.1980/}{CoDial: Interpretable Task-Oriented Dialogue Systems Through Dialogue Flow Alignment}}
  {\textbf{Radin Shayanfar}, Chu Fei Luo, Rohan V.\ Bhambhoria, Samuel Dahan, Xiaodan Zhu.}
  {ACL 2026}

\pubentry[Automated red teaming in which multiple LLMs debate to surface each other's unsafe behaviour, with long-term memory carrying safety lessons across debates.]
  {\href{https://openreview.net/forum?id=79fSrxFKKx}{RedDebate: Safer Responses through Multi-Agent Red Teaming Debates}}
  {Ali Asad, Stephen Obadinma, \textbf{Radin Shayanfar}, Xiaodan Zhu.}
  {ICML 2026}

\pubentry
  {\href{https://centre-for-software-excellence.github.io/docs/blog/before-you-score-the-model-score-the-benchmark}{Before You Score the Model, Score the Benchmark: A Skeptical View Into Current Agentic Software Engineering Benchmarks}}
  {\textbf{Radin Shayanfar}, Keheliya Gallaba.}
  {Blog Post\\May 2026}

\pubentry
  {\href{https://aclanthology.org/2024.findings-emnlp.924/}{Misinformation with Legal Consequences (MisLC): A New Task Towards Harnessing Societal Harm of Misinformation}}
  {Chu Fei Luo, \textbf{Radin Shayanfar}, Rohan V.\ Bhambhoria, Samuel Dahan, Xiaodan Zhu.}
  {EMNLP 2024\\(Findings)}

\section*{Open-Source Contributions}

\entry{\href{https://github.com/harbor-framework/harbor}{Harbor Framework}}{Agent-evaluation harness and RL-environment framework from the creators of Terminal-Bench}{Contributor}{2026}
\begin{itemize}
  \item Authored the FeatBench adapter for Harbor (\href{https://github.com/harbor-framework/harbor/pull/1218}{PR \#1218}), a 156-task, 25-repo feature-level code-generation benchmark, and ported the parity fixes upstream, preserving 99\%+ of original instances and tests.
  \item Integrated ByteDance's trae-agent into Harbor (\href{https://github.com/harbor-framework/harbor/pull/1298}{PR \#1298}) with full agent implementation, configuration handling, API-key sanitization, and unit tests.
  \item Fixed an environment-variable bug in deveval oracle script generation (\href{https://github.com/harbor-framework/harbor/pull/1431}{PR \#1431}).
  \item Reviewed the SWE-Gym adapter fix for pytest-xdist compatibility (\href{https://github.com/harbor-framework/harbor/pull/1438}{PR \#1438}); flagged that the proposed universal flag would break environments without pytest-xdist installed.
\end{itemize}

\entry{\href{https://github.com/NVIDIA-NeMo/Guardrails}{NVIDIA NeMo Guardrails}}{Toolkit for programmable guardrails on LLM-based conversational systems}{Contributor}{2024}
\begin{itemize}
  \item Fixed the Colang~2 parser to correctly handle nested \texttt{if/else} statements (\href{https://github.com/NVIDIA-NeMo/Guardrails/pull/833}{PR \#833}).
  \item Resolved multiline output rendering in few-shot prompts (\href{https://github.com/NVIDIA-NeMo/Guardrails/pull/765}{PR \#765}) and dynamic flow generation (\href{https://github.com/NVIDIA-NeMo/Guardrails/pull/748}{PR \#748}).
  \item Added Colang commenting support (\texttt{Ctrl+/}) to the VSCode extension (\href{https://github.com/NVIDIA-NeMo/Guardrails/pull/749}{PR \#749}).
\end{itemize}

\section*{Projects}

\entry{\href{https://github.com/radinshayanfar/transformer}{Transformer from Scratch}}{PyTorch, Transformer, Machine Translation, BPE Tokenization}{}{Summer 2025}
\begin{itemize}
  \item Implemented the original Transformer architecture in PyTorch with encoder-decoder, encoder-only, and decoder-only variants --- working solely from the original paper, without referencing external implementations.
  \item Built an end-to-end translation system on the WMT dataset: data loading, BPE tokenization, batching, masking, training loop, and evaluation.
  \item Structured the codebase to extend cleanly to text generation and Vision Transformer (ViT) tasks on the same base.
\end{itemize}

\entry{\href{https://github.com/radinshayanfar/speechbrain}{Speaker-Specific Voice Activity Detection with ECAPA-TDNN Embeddings}}{PyTorch, SpeechBrain, VAD, Speaker Verification}{}{Fall 2024}
\begin{itemize}
  \item Proposed and implemented a Speaker-Specific VAD (SVAD) framework integrating ECAPA-TDNN embeddings into a CRDNN-based VAD.
  \item Extended SpeechBrain to support conditioning on target-speaker embeddings by modifying model internals and data pipelines.
  \item Engineered training and augmentation pipelines (including embedding-level noise injection) to improve generalization under speaker-limited data.
  \item Evaluated on the AMI dataset with low Diarization Error Rate under real-time, multi-speaker conditions.
\end{itemize}

\entry{\href{https://github.com/aghassel/NLP-Sentiment-Analysis/}{Comparative Analysis of Transformer-Based Models for Sentiment Analysis}}{PyTorch, Transformers, Sentiment Analysis}{}{Fall 2023}
\begin{itemize}
  \item Fine-tuned GPT-2 for multi-label emotion and binary sentiment classification on the GoEmotions and Sent140 datasets.
  \item Added a classification head and applied dropout-based regularization, showing GPT-2 is competitive with BERT and ELMo on these tasks.
\end{itemize}

\entry{Deep Recurrent Q-Network Agent for a 3D FPS Game}{PyTorch, Reinforcement Learning, DRQN}{}{Fall 2023}
\begin{itemize}
  \item Reimplemented the \href{https://arxiv.org/pdf/1609.05521}{DRQN} from scratch to train an autonomous VizDoom agent from visual observations.
  \item Integrated LSTM layers for partial observability and added game-feature prediction as an auxiliary task.
\end{itemize}

\entry{Generating Handwritten Persian Characters with Generative Models}{Python, Keras, PyTorch, Generative Models}{}{Feb.\ 2022 -- Apr.\ 2023}
\begin{itemize}
  \item Studied the ability of GAN (and DCGAN, WGAN-GP), VAE, Real NVP, and DDPM to represent individual Persian handwritten characters.
  \item Trained a Real NVP model and developed a novel technique to replicate personal handwriting styles from sample data. B.Sc.\ project under Prof.\ A.\ Nickabadi.
\end{itemize}

\entry{\href{https://github.com/radinshayanfar/AUT-IR}{Information Retrieval Engine}}{Python}{}{2022}
\begin{itemize}
  \item Built a text-based search engine with Boolean and TF-IDF ranked retrieval, using a positional inverted index.
\end{itemize}

\entry{Habco --- Healthcare Backend}{Laravel, MySQL, Redis, REST API}{}{2021}
\begin{itemize}
  \item Built a RESTful API in Laravel for a Canadian healthcare startup serving patients, doctors, nurses, and pharmacists --- prescription sharing, drug inventory, SMS-based authentication.
\end{itemize}

\entry{Music Identification}{Python, NumPy, Signals and Systems}{}{Spring 2021}
\begin{itemize}
  \item Designed and implemented an audio-fingerprinting and music-identification system from scratch using FFT and spectrogram analysis with a custom similarity metric.
  \item Assigned this project as the final assignment for the undergraduate Signals and Systems course while serving as teaching assistant.
\end{itemize}

\entry{\href{https://samcode.allamehelli3.ir/}{SamCode Website}}{PHP, HTML, CSS, Bootstrap, Sass, JavaScript}{}{2020}
\begin{itemize}
  \item Designed and developed the official website for SamCode, a programming contest hosted by Allameh Helli 3 Junior High School.
  \item Customized the Datalife Engine (DLE) CMS with a new theme integrating a modern landing page with the system's default template.
\end{itemize}

\entry{Captcha Solver}{Python, Keras, OpenCV}{}{2019 -- 2020}
\begin{itemize}
  \item Built a CAPTCHA-solving system with OpenCV preprocessing and a Keras neural network, achieving 96\% accuracy on 5-character alphanumeric CAPTCHAs.
  \item Collected and annotated training data via a custom web-based annotation UI.
\end{itemize}

\entry{\href{https://github.com/radinshayanfar/RJBot}{RJBot}}{PHP, MySQL, Telegram Bot API}{}{2019 -- 2020}
\begin{itemize}
  \item Developed an open-source Telegram bot for searching and downloading media from RadioJavan.com, supporting music, videos, albums, and podcasts.
\end{itemize}

\entry{\href{https://github.com/radinshayanfar/Jpotify}{Jpotify}}{Java, Swing}{}{2019}
\begin{itemize}
  \item Created a desktop music player in Java using Swing for the graphical interface, with playlist management and music sharing over a local network.
\end{itemize}

\entry{Ping-Pong Ball Tracker}{C++, OpenCV}{}{2014}
\begin{itemize}
  \item Developed a real-time ping-pong ball tracking system using OpenCV and C++.
\end{itemize}

\section*{Skills}
\begin{itemize}
  \item \textbf{Programming:} Python, Go, Java, C, JavaScript, LaTeX, LISP, ML.
  \item \textbf{Generative AI \& Agents:} large language models, agentic AI, coding agents, agent scaffolds, agent evaluation and benchmarking, benchmark construction, trajectory analysis, sandboxed execution, tool use, MCP, SWE-Bench, OpenHands, Harbor, Terminal-Bench, vLLM, LangChain, LiteLLM, NeMo Guardrails, RAG, Pyserini, BM25, Pinecone.
  \item \textbf{Machine Learning \& NLP:} PyTorch, HuggingFace Transformers, Keras, SpeechBrain, data science, statistical analysis.
  \item \textbf{Big Data \& Infra:} Apache Spark, Apache Hive, Apache Hadoop, Apache Kafka, Docker, Kubernetes, Linux, Git, server administration.
  \item \textbf{Other:} CUDA, OpenMP, OpenCV, Arduino.
\end{itemize}

\section*{Certificates}
\begin{itemize}
  \item \textbf{AWS Technical Essentials} --- \href{https://explore.skillbuilder.aws/learn/courses/1851/aws-technical-essentials}{Amazon Web Services}, Spring 2025.
  \item \textbf{\href{https://www.coursera.org/account/accomplishments/specialization/certificate/WT5VCNDDCB7V}{Generative Adversarial Networks Specialization}} --- DeepLearning.AI, Spring 2022.
  \item \textbf{\href{https://www.coursera.org/account/accomplishments/certificate/ZVWNNHWAJUWC}{Game Theory}} --- Stanford University \& The University of British Columbia, Summer 2021.
  \item \textbf{\href{https://www.coursera.org/account/accomplishments/specialization/certificate/3DWHZC46YSCK}{Deep Learning Specialization}} --- DeepLearning.AI, Summer--Fall 2020.
  \item \textbf{\href{https://www.coursera.org/account/accomplishments/certificate/2XNJSGPREAQY}{Machine Learning}} --- Stanford University, Spring 2020.
  \item \textbf{Reinforcement Learning} --- University College London (David Silver, audited), Summer 2020.
\end{itemize}

\section*{Honors \& Awards}
\begin{itemize}
  \item Top 0.3\% among 130{,}000+ applicants, \href{https://en.wikipedia.org/wiki/Iranian_University_Entrance_Exam}{Nationwide University Entrance Exam} (B.Sc.\ in Mathematics and Engineering), Iran, 2018.
  \item 3rd place, Iran Zamin Open Cup (IZOCup) Programming Contest, Tehran, 2014.
  \item Qualified, National Organization for Development of Exceptional Talents (NODET) entrance exam, Iran, 2011.
\end{itemize}

\section*{Languages}
\begin{itemize}
  \item \textbf{Persian:} Native.
  \item \textbf{English:} Proficient. TOEFL (2022): 109 / 120 (R: 29, L: 30, S: 21, W: 29).
\end{itemize}

\section*{References}

\reference{Dr.\ Xiaodan Zhu}{Associate Professor, ECE}{Queen's University}{xiaodan.zhu@queensu.ca}{https://scholar.google.ca/citations?user=a6MYnuUAAAAJ}

\reference[dblp]{Dr.\ Thomas Dean}{Professor, ECE}{Queen's University}{tom.dean@queensu.ca}{https://dblp.org/pid/d/ThomasRDean.html}

\reference{Dr.\ Michael Greenspan}{Professor, ECE}{Queen's University}{greenspan.michael@queensu.ca}{https://scholar.google.ca/citations?user=x-M8l7IAAAAJ}

\reference{Dr.\ Seyed Abolfazl Motahari}{Assistant Professor, CE}{Sharif University of Technology}{motahari@sharif.edu}{https://scholar.google.com/citations?user=rJ-biB0AAAAJ}

\reference{Dr.\ Ahmad Nickabadi}{Assistant Professor, CE}{Amirkabir University of Technology}{nickabadi@aut.ac.ir}{https://scholar.google.com/citations?user=pSMNSZwAAAAJ}

\reference{Dr.\ Mehdi Rasti}{Associate Professor, ITEE}{University of Oulu}{mehdi.rasti@oulu.fi}{https://scholar.google.com/citations?user=zb8pjMYAAAAJ}

\end{document}
